\documentclass[11pt]{article}
\usepackage[margin=0.4in]{geometry}
\usepackage{graphicx}
\usepackage{amsmath}
\usepackage{booktabs}
\usepackage{microtype}
\usepackage[colorlinks=true, linkcolor=blue, urlcolor=blue, citecolor=blue]{hyperref}

% Compact title formatting
\usepackage{titling}
\pretitle{\begin{center}\large\bfseries}
\posttitle{\end{center}\vspace{-1.5em}}
\preauthor{\begin{center}\normalsize}
\postauthor{\end{center}\vspace{-1.5em}}
\predate{\begin{center}\normalsize}
\postdate{\end{center}\vspace{-2em}}

\title{Project Status Update: Electrolyte Optimization\\for Sorbent Atmospheric Water Harvesting}
\author{Veronica Fish and Clara Drysdale}
\date{May 15, 2026}

\begin{document}
\maketitle
\vspace{-2em}

\section*{Introduction}
\vspace{-0.8em}

Atmospheric water harvesting devices pull drinking water directly from humid air. The devices we study use a salt-infused polymer gel that passively absorbs moisture overnight, then releases liquid water when heated by sunlight during the day. The salt blend and the ratio of salt to polymer strongly determine how much water is produced and at what cost. We are building an optimization pipeline to find the best formulation for any location on Earth, minimizing the levelized cost of water (LCOW, \$/kg).
\vspace{-1em}
\section*{Optimization Formulation}
\vspace{-0.8em}
\textbf{How the physics feeds the optimizer.}
The amount of water a given salt blend can capture depends on local humidity and temperature: each salt absorbs water up to a characteristic equilibrium, and the difference between the overnight absorption state and the daytime release state determines the daily water yield. We compute these equilibria numerically from empirical salt data before running the optimizer.

\textbf{What is being optimized.}
Decision variables are blend weights $f_i\!\in\![0,1]$ for four salts (LiCl, NaCl, CaCl$_2$, MgCl$_2$) and salt-to-polymer ratio $s/l\!\in\![0.05,16]$. The ZSR rule gives blended molality $m_i = f_i\,m_i^\circ(a_w,T)$; the uptake coefficient $B = \frac{x_w\nu}{1-x_w}\frac{M_w}{\bar{M}_s}$ at each operating point yields the half-swing $\Delta B = B_\mathrm{up}-B_\mathrm{des}$. The NLP minimizes
\begin{equation*}
  \mathrm{LCOW} = \frac{\bigl(\mathrm{CRF}+f_m\bigr)C_\mathrm{dev} + C_\mathrm{gel}/\tau}
                       {\eta\cdot 365\cdot\tfrac{s/l}{1+s/l}\cdot\Delta B\cdot m_\mathrm{dry}}
\end{equation*}
subject to $\sum_i f_i = 1$ and $\Delta B \geq 0$, where $\mathrm{CRF}=i(1+i)^L/[(1+i)^L-1]$. Ipopt (interior-point) solves the NLP; a second mode traces the LCOW--yield Pareto front via epsilon-constraint scalarization.
\vspace{-1em}
\section*{Results}
\vspace{-0.8em}
\begin{figure}[h]
  \centering
  \includegraphics[width=0.5\linewidth]{lcow_random_sites.png}
  \caption{Optimized LCOW (\$/kg water) at randomly sampled land sites; color encodes cost and marker shape encodes the dominant salt in the optimal blend.}
  \label{fig:results}
\end{figure}

Of 60 randomly sampled land sites, 52 are feasible---the other 8 are too dry or cold for any salt blend to complete a full absorption--release cycle. Lithium chloride (LiCl) is the dominant salt at 47 of 52 sites because it remains effective across a wide range of humidity levels; magnesium chloride takes over at 4 cold Arctic sites. Optimized LCOW ranges from \$0.002 to \$0.090/kg (median \$0.009/kg), competitive with or below municipal water costs (Figure~\ref{fig:results}).
\vspace{-1em}
\section*{Next Steps}
\vspace{-0.8em}
The pipeline is fully implemented and unit-tested. Priorities are: (1) extending the model to active SAWH devices; (2) sensitivity analyses on key economic assumptions---device lifetime, material replacement intervals, and capital cost---to understand which parameters most drive LCOW uncertainty; and (3) adding a toxicity penalty to the objective.

\end{document}